%%%%%%%% ICML 2026 EXAMPLE LATEX SUBMISSION FILE %%%%%%%%%%%%%%%%%

\documentclass{article}

% Recommended, but optional, packages for figures and better typesetting:
\usepackage{microtype}
\usepackage{graphicx}
\usepackage{subcaption}
\usepackage{booktabs} % for professional tables

% hyperref makes hyperlinks in the resulting PDF.
% If your build breaks (sometimes temporarily if a hyperlink spans a page)
% please comment out the following usepackage line and replace
% \usepackage{icml2026} with \usepackage[nohyperref]{icml2026} above.
\usepackage{hyperref}


% Attempt to make hyperref and algorithmic work together better:
\newcommand{\theHalgorithm}{\arabic{algorithm}}

% Use the following line for the initial blind version submitted for review:
\usepackage{icml2026}

% For preprint, use
% \usepackage[preprint]{icml2026}

% If accepted, instead use the following line for the camera-ready submission:
% \usepackage[accepted]{icml2026}

\usepackage{amsmath}
\usepackage{amssymb}
\usepackage{mathtools}
\usepackage{amsthm}


% if you use cleveref..
\usepackage[capitalize,noabbrev]{cleveref}

%%%%%%%%%%%%%%%%%%%%%%%%%%%%%%%%
% THEOREMS
%%%%%%%%%%%%%%%%%%%%%%%%%%%%%%%%
\theoremstyle{plain}
\newtheorem{theorem}{Theorem}[section]
\newtheorem{proposition}[theorem]{Proposition}
\newtheorem{lemma}[theorem]{Lemma}
\newtheorem{corollary}[theorem]{Corollary}
\theoremstyle{definition}
\newtheorem{definition}[theorem]{Definition}
\newtheorem{assumption}[theorem]{Assumption}
\theoremstyle{remark}
\newtheorem{remark}[theorem]{Remark}

% Todonotes is useful during development; simply uncomment the next line
%    and comment out the line below the next line to turn off comments
%\usepackage[disable,textsize=tiny]{todonotes}
\usepackage[textsize=tiny]{todonotes}

% The \icmltitle you define below is probably too long as a header.
% Therefore, a short form for the running title is supplied here:
\icmltitlerunning{Submission and Formatting Instructions for ICML 2026}

\begin{document}

\twocolumn[
  \icmltitle{Submission and Formatting Instructions for \\
    International Conference on Machine Learning (ICML 2026)}

  % It is OKAY to include author information, even for blind submissions: the
  % style file will automatically remove it for you unless you've provided
  % the [accepted] option to the icml2026 package.

  % List of affiliations: The first argument should be a (short) identifier you
  % will use later to specify author affiliations Academic affiliations
  % should list Department, University, City, Region, Country Industry
  % affiliations should list Company, City, Region, Country

  % You can specify symbols, otherwise they are numbered in order. Ideally, you
  % should not use this facility. Affiliations will be numbered in order of
  % appearance and this is the preferred way.
  \icmlsetsymbol{equal}{*}

  \begin{icmlauthorlist}
    \icmlauthor{Firstname1 Lastname1}{equal,yyy}
    \icmlauthor{Firstname2 Lastname2}{equal,yyy,comp}
    \icmlauthor{Firstname3 Lastname3}{comp}
    \icmlauthor{Firstname4 Lastname4}{sch}
    \icmlauthor{Firstname5 Lastname5}{yyy}
    \icmlauthor{Firstname6 Lastname6}{sch,yyy,comp}
    \icmlauthor{Firstname7 Lastname7}{comp}
    %\icmlauthor{}{sch}
    \icmlauthor{Firstname8 Lastname8}{sch}
    \icmlauthor{Firstname8 Lastname8}{yyy,comp}
    %\icmlauthor{}{sch}
    %\icmlauthor{}{sch}
  \end{icmlauthorlist}

  \icmlaffiliation{yyy}{Department of XXX, University of YYY, Location, Country}
  \icmlaffiliation{comp}{Company Name, Location, Country}
  \icmlaffiliation{sch}{School of ZZZ, Institute of WWW, Location, Country}

  \icmlcorrespondingauthor{Firstname1 Lastname1}{first1.last1@xxx.edu}
  \icmlcorrespondingauthor{Firstname2 Lastname2}{first2.last2@www.uk}

  % You may provide any keywords that you find helpful for describing your
  % paper; these are used to populate the "keywords" metadata in the PDF but
  % will not be shown in the document
  \icmlkeywords{Machine Learning, ICML}

  \vskip 0.3in
]

% this must go after the closing bracket ] following \twocolumn[ ...

% This command actually creates the footnote in the first column listing the
% affiliations and the copyright notice. The command takes one argument, which
% is text to display at the start of the footnote. The \icmlEqualContribution
% command is standard text for equal contribution. Remove it (just {}) if you
% do not need this facility.

% Use ONE of the following lines. DO NOT remove the command.
% If you have no special notice, KEEP empty braces:
\printAffiliationsAndNotice{}  % no special notice (required even if empty)
% Or, if applicable, use the standard equal contribution text:
% \printAffiliationsAndNotice{\icmlEqualContribution}

\begin{abstract}
  This document provides a basic paper template and submission guidelines.
  Abstracts must be a single paragraph, ideally between 4--6 sentences long.
  Gross violations will trigger corrections at the camera-ready phase.
\end{abstract}

\section{Introduction}

The rapid adoption of modern large language model (LLM) inference engines has created a critical need for optimization benchmarks that reflect real-world production scenarios. However,  existing benchmarks fail to assess real-world optimization scenarios, leading to a gap between research and production needs. We introduce the Inference Optimization Benchmark (IOBench), a methodology for evaluating LLM-driven code optimization across diverse GPU hardware. IOBench employs a three-stage pipeline that analyzes repository commit histories from production inference engines (VLLM, SGLang, Max Inference Engine), identifies and distills performance-critical commits through automated analysis and manual refinement (~1,700-3,000 commits reduced to 100+), and extracts minimal reproducible optimization tasks with documented performance characteristics. 
The Benchmark tasks span attention mechanisms, memory management, quantization, and scheduling, with difficulty levels designed to reflect realistic engineering complexity. Our evaluation framework measures correctness, performance gains, hardware efficiency, and cross-platform portability on NVIDIA (A100, H100) and AMD (MI300X) GPUs. By bridging the gap between code-level optimization and large-scale deployment, IOBench provides a reproducible and extensible benchmark that advances the study of inference optimization in LLM systems.


\section{Related Work}
\label{sec:related}

Our work connects four strands of prior research: (i) benchmarks for code generation, (ii) performance-aware code generation and optimization, (iii) agent architectures for software engineering, and (iv) repository mining methodologies. We briefly review each area and highlight how \emph{IOBench} (OmniPerf-Bench) differs.

\paragraph{Benchmarks for code generation.}
Large language model evaluation has focused primarily on functional correctness.
\textbf{HumanEval}~\citep{chen2021evaluating} introduced 164 function-level programming problems with the pass@k metric.
\textbf{SWE-bench}~\citep{jimenez2023swe} scaled to 2,294 real-world GitHub issues requiring repository-level edits, becoming the standard for evaluating coding agents.

However, these benchmarks measure only whether code \emph{works}, not whether it performs \emph{well}---a critical gap for production systems where performance directly impacts deployment cost and latency.

\paragraph{Performance-aware code generation and optimization.}
Recent work has begun addressing performance optimization.
\textbf{KernelBench}~\citep{ouyang2025kernelbench} evaluates LLMs on GPU kernel generation, finding that frontier models match baseline performance in less than 20\% of cases.
\textbf{TritonBench}~\citep{li2025tritonbench} focuses on Triton operator generation, revealing gaps in producing efficient low-level code.

Moving to general software optimization, two recent benchmarks are most relevant to our work.
\textbf{GSO}~\citep{shetty2025gso} curates 102 optimization tasks across 10 repositories using commit history mining, reporting <5\% agent success and identifying failures in bottleneck localization.
\textbf{SWE-Perf}~\citep{he2025sweperf} provides 140 instances from performance-enhancing GitHub PRs, evaluating file-level and repository-level approaches.

While GSO and SWE-Perf establish baseline difficulty, they report primarily aggregate success rates. IOBench extends this by providing \emph{behavioral analysis}---identifying specific failure modes (instant-edit pathology predicting 95\% failure, commit explosion up to 7,755 commits) and codebase complexity effects (58-point performance variance using identical agent configuration).

\paragraph{Agent architectures for software engineering.}
\textbf{OpenHands}~\citep{wang2024openhands} and \textbf{SWE-Agent}~\citep{yang2024sweagent} are representative open-source frameworks for coding agents, demonstrating that agent-computer interface design significantly impacts performance on SWE-bench.

Comparative evaluation across agent architectures on identical tasks is rare. IOBench evaluates three distinct architectures (Codex, TRAE, OpenHands) on 179 optimization tasks from production ML inference engines (99 from vLLM, 80 from SGLang), revealing architectural trade-offs in cost, success rate, and observability.

\paragraph{Repository mining for benchmarks.}
Our methodology draws on prior work mining software repositories for realistic evaluation datasets.
\textbf{BugSwarm}~\citep{tomassi2019bugswarm} mines CI logs to create reproducible failure-fix pairs, demonstrating that automated repository mining scales better than manual curation.
GSO and SWE-Perf (discussed above) apply similar techniques to performance optimization.

IOBench mines commit histories from production ML inference engines (\textbf{vLLM}~\citep{kwon2023vllm} with PagedAttention for KV-cache management, and \textbf{SGLang}~\citep{zheng2023sglang} for structured generation).
These systems represent heavily-optimized production codebases whose performance results from incremental engineering changes.
By extracting optimization tasks from their evolution, IOBench provides domain-relevant challenges for ML systems applications.

\section{Methodology}
\label{sec:methodology}

\todo{To be written: Dataset construction pipeline, commit mining, test generation, evaluation framework}

\section{Experimental Setup}
\label{sec:experiments}

\todo{To be written: Agent configurations, hardware setup, evaluation metrics}

\section{Results}
\label{sec:results}

\todo{To be written: Success rates, behavioral patterns, codebase effects}

\section{Analysis and Discussion}
\label{sec:analysis}

\todo{To be written: Failure mode analysis, instant-edit pathology, commit explosion, violations}

\section{Conclusion}
\label{sec:conclusion}

\todo{To be written: Summary, implications, future work}

\section*{Acknowledgments}

\todo{To be added}

% Bibliography
\bibliography{references}
\bibliographystyle{icml2026}

\end{document}
